\documentclass[11pt,a4paper]{article}
\usepackage[top=2.3cm,bottom=2.3cm,left=2.2cm,right=2.2cm]{geometry}

\usepackage{amsmath,amssymb}
\usepackage{fontspec}
\setmainfont{Liberation Serif}
\setsansfont{Liberation Sans}
\setmonofont{DejaVu Sans Mono}[Scale=0.84]

\usepackage{xcolor}
\definecolor{brandblue}{RGB}{20, 55, 105}
\definecolor{accentblue}{RGB}{35, 95, 165}
\definecolor{codebg}{RGB}{248, 249, 251}
\definecolor{codeframe}{RGB}{215, 222, 230}
\definecolor{javakeyword}{RGB}{0, 45, 170}
\definecolor{javastring}{RGB}{5, 120, 20}
\definecolor{javacomment}{RGB}{120, 120, 120}
\definecolor{javanumber}{RGB}{20, 75, 220}

\definecolor{noteborder}{RGB}{30, 110, 65}
\definecolor{notebg}{RGB}{242, 249, 245}
\definecolor{warnborder}{RGB}{185, 75, 10}
\definecolor{warnbg}{RGB}{254, 248, 240}
\definecolor{tipborder}{RGB}{30, 95, 160}
\definecolor{tipbg}{RGB}{242, 247, 254}

\usepackage{fancyhdr}
\usepackage{lastpage}
\pagestyle{fancy}
\fancyhf{}
\lhead{\parbox[t]{0.58\textwidth}{\raggedright\small\sffamily\color{brandblue}\textbf{T02 --- Paradigma ad Oggetti e Modellazione UML}}}
\rhead{\small\sffamily\color{gray}Programmazione II -- UniVR (2025/2026)}
\cfoot{\small\sffamily Pagina \thepage\ di \pageref{LastPage}}
\renewcommand{\headrulewidth}{0.5pt}
\renewcommand{\headrule}{\hbox to\headwidth{\color{brandblue!70!black}\leaders\hrule height \headrulewidth\hfill}}

\usepackage{titlesec}
\titleformat{\section}{\Large\bfseries\sffamily\color{brandblue}}{\thesection}{0.8em}{}[\titlerule]
\titleformat{\subsection}{\large\bfseries\sffamily\color{accentblue}}{\thesubsection}{0.8em}{}
\titleformat{\subsubsection}{\normalsize\bfseries\sffamily\color{brandblue!85!black}}{\thesubsubsection}{0.8em}{}

\usepackage{needspace}
\usepackage{fancyvrb}
\DefineVerbatimEnvironment{asciibox}{Verbatim}{
  fontsize=\fontsize{7.5pt}{9.2pt}\selectfont,
  frame=single,
  rulecolor=\color{codeframe},
  fillcolor=\color{codebg},
  framesep=2.5mm
}
\usepackage{listings}
\lstset{
  backgroundcolor=\color{codebg},
  frame=single,
  rulecolor=\color{codeframe},
  basicstyle=\ttfamily\footnotesize,
  keywordstyle=\color{javakeyword}\bfseries,
  stringstyle=\color{javastring},
  commentstyle=\color{javacomment}\itshape,
  numberstyle=\tiny\color{gray},
  numbers=left,
  numbersep=7pt,
  stepnumber=1,
  breaklines=true,
  breakatwhitespace=true,
  showstringspaces=false,
  tabsize=4,
  xleftmargin=12pt,
  framexleftmargin=12pt
}

\newcommand{\passthrough}[1]{#1}
\providecommand{\tightlist}{\setlength{\itemsep}{0pt}\setlength{\parskip}{0pt}}

\usepackage{tcolorbox}
\tcbuselibrary{skins,breakable}
\newtcolorbox{studynote}[1][Nota]{
  enhanced,
  breakable,
  colback=notebg,
  colframe=noteborder,
  fonttitle=\bfseries\sffamily\small,
  coltitle=white,
  title={#1},
  arc=2mm,
  boxrule=0.8pt,
  left=9pt, right=9pt, top=7pt, bottom=7pt
}

\newtcolorbox{studywarning}[1][Attenzione]{
  enhanced,
  breakable,
  colback=warnbg,
  colframe=warnborder,
  fonttitle=\bfseries\sffamily\small,
  coltitle=white,
  title={#1},
  arc=2mm,
  boxrule=0.8pt,
  left=9pt, right=9pt, top=7pt, bottom=7pt
}

\newtcolorbox{studytip}[1][Suggerimento]{
  enhanced,
  breakable,
  colback=tipbg,
  colframe=tipborder,
  fonttitle=\bfseries\sffamily\small,
  coltitle=white,
  title={#1},
  arc=2mm,
  boxrule=0.8pt,
  left=9pt, right=9pt, top=7pt, bottom=7pt
}

\usepackage{booktabs}
\usepackage{longtable}
\usepackage{array}
\usepackage{calc}

\usepackage{hyperref}
\hypersetup{
  colorlinks=true,
  linkcolor=accentblue,
  urlcolor=accentblue,
  citecolor=brandblue,
  pdfborder={0 0 0}
}

\title{\Huge\bfseries\sffamily\color{brandblue} T02 --- Paradigma ad Oggetti e Modellazione UML \\[0.3em] \large\sffamily Astrazione, Incapsulamento, Modularità, Gerarchia e Sintassi dei Class Diagram}
\author{\textbf{Docente:} Prof. Michele Pasqua \quad---\quad \textbf{Appunti Revisione:} Wally}
\date{A.A. 2025/2026 -- Universit\`a degli Studi di Verona}

\begin{document}
\maketitle
\thispagestyle{fancy}
\vspace{-1.2em}
\noindent\textcolor{brandblue}{\rule{\textwidth}{0.8pt}}
\vspace{1.2em}

In questo blocco analizziamo la transizione teorico-metodologica dal
paradigma procedurale a quello ad oggetti, i concetti cardine dell'OOP e
la modellazione grafica tramite \textbf{UML}
(\href{file:///home/wally/Documenti/GitHub/Programmazione-II/25-26/Teoria-e-Esercitazioni/Slides-20260902/T02\%20-\%20Object-Orientations\%20and\%20UML.pdf}{T02
- Object-Orientations and UML.pdf}).

\begin{center}\rule{0.5\linewidth}{0.5pt}\end{center}

\subsubsection{1. Teoria: Concetti Teorici dalle
Slide}\label{teoria-concetti-teorici-dalle-slide}

\paragraph{A. La Tassonomia dei Linguaggi e il Bisogno di Astrazione
(Slide
1--13)}\label{a.-la-tassonomia-dei-linguaggi-e-il-bisogno-di-astrazione-slide-113}

\begin{itemize}
\tightlist
\item
  \textbf{Nessun linguaggio universale}: citando Edsger W. Dijkstra
  (\emph{«Computer Science is no more about computers than Astronomy is
  about telescopes»}), il docente introduce il principio del
  \textbf{Golden Hammer} (\emph{«If all you have is a hammer, everything
  looks like a nail»}): non esiste un linguaggio perfetto per tutto.
\item
  \textbf{Evoluzione storica degli obiettivi}:

  \begin{itemize}
  \tightlist
  \item
    \emph{Anni '50/'60}: focus sull'\textbf{efficienza del codice
    macchina}. L'hardware costava milioni e i programmatori poco:
    l'obiettivo era ``tenere occupata la macchina'' con linguaggi vicini
    al silicio.
  \item
    \emph{Oggi}: focus sull'\textbf{efficienza dello sviluppo del
    software}. L'hardware costa poco e i programmatori costano molto:
    l'obiettivo è ``tenere produttivo il programmatore'', abbattendo
    complessità, bug e costi di manutenzione tramite
    l'\textbf{astrazione}.
  \end{itemize}
\item
  \textbf{I Tre Paradigmi Fondamentali}:

  \begin{enumerate}
  \def\labelenumi{\arabic{enumi}.}
  \tightlist
  \item
    \textbf{Imperativo / Procedurale} (C, Pascal, Fortran, COBOL):

    \begin{itemize}
    \tightlist
    \item
      Basato sul concetto di \textbf{stato} (celle di memoria che
      memorizzano valori) e \textbf{comandi sequenziali} che mutano tale
      stato (assegnamenti).
    \item
      Riuso del codice tramite procedure/funzioni che operano su dati
      esterni.
    \end{itemize}
  \item
    \textbf{Dichiarativo}:

    \begin{itemize}
    \tightlist
    \item
      \emph{Funzionale} (Haskell, Lisp): il programma è la valutazione
      di un'espressione matematica pura, basata su applicazione di
      funzioni e ricorsione (es.
      \passthrough{\lstinline!add x y = x + y!}). Assenza o
      minimizzazione degli effetti collaterali (\emph{side effects}).
    \item
      \emph{Logico} (Prolog): il programma definisce relazioni, fatti e
      regole di deduzione logica; l'esecuzione è una
      risoluzione/unificazione guidata da query (es.
      \passthrough{\lstinline!playsAirGuitar(sam) :- listens2Music(sam).!}).
    \end{itemize}
  \item
    \textbf{Object-Oriented (OOP)} (Smalltalk, Java, C++):

    \begin{itemize}
    \tightlist
    \item
      Unisce il nucleo imperativo (gli oggetti contengono uno stato
      interno) a tratti dichiarativi (gli oggetti espongono
      comportamenti invocabili senza conoscerne l'algoritmo interno).
    \item
      Basato su \textbf{istanze} che comunicano tramite
      \textbf{passaggio di messaggi}.
    \end{itemize}
  \end{enumerate}
\item
  \textbf{Lo Spettro dell'Astrazione e il Trade-off di Overhead}:
  \[\text{Machine Code} \longrightarrow \text{Assembly} \longrightarrow \text{Procedurale} \longrightarrow \text{Object-Oriented} \longrightarrow \text{Funzionale}\]

  \begin{itemize}
  \tightlist
  \item
    Più si sale nell'astrazione:

    \begin{itemize}
    \tightlist
    \item
      \textbf{Programmer Overhead \(\searrow\) (minimo)}: meno tempo
      perso a gestire registri, puntatori e memoria manuale.
    \item
      \textbf{Compiler/Runtime Overhead \(\nearrow\) (massimo)}: il
      compilatore/JVM deve effettuare controlli di tipo a runtime,
      garbage collection e risoluzione dinamica delle chiamate
      (\emph{dynamic dispatch}).
    \end{itemize}
  \end{itemize}
\item
  \textbf{Linea temporale}: dal 1949 (Assembly), Fortran (1957), C
  (1972), C++ (1983) fino all'avvento di \textbf{Java (1995)}, C\#
  (2000) e linguaggi moderni multiparadigma (Swift 2014).
\end{itemize}

\begin{center}\rule{0.5\linewidth}{0.5pt}\end{center}

\paragraph{B. La Crisi del Procedurale e la Nascita degli Oggetti (Slide
14--23)}\label{b.-la-crisi-del-procedurale-e-la-nascita-degli-oggetti-slide-1423}

Il prof. Pasqua presenta il caso di studio di un programma in C che
manipola un vettore globale:

\begin{lstlisting}[language=C]
int vect[20];
void sort() { /* ordina vect */ }
int search(int n) { /* cerca in vect */ }
void init() { /* inizializza vect */ }
int i;

void main() {
    init();
    sort();
    search(13);
}
\end{lstlisting}

\textbf{I 4 Difetti Letali dell'Approccio Procedurale evidenziati a
lezione}: 1. \textbf{Relazione Implicita Dati-Funzioni}:
\passthrough{\lstinline!vect!} e le funzioni
\passthrough{\lstinline!sort()!}, \passthrough{\lstinline!search()!}
sono entità slegate a livello di linguaggio. Non esiste un tipo autonomo
``Vettore'' che protegga se stesso. 2. \textbf{Nessun Controllo dei
Confini (Out-of-Bounds)}: un ciclo come
\passthrough{\lstinline!for (i=0; i<=21; i++) vect[i]=0;!} scrive in
memoria oltre i 20 elementi, corrompendo altre variabili o lo stack
senza alcun errore del compilatore (\emph{buffer overflow}). 3.
\textbf{Nessuna Garanzia sull'Inizializzazione}: se il programmatore
dimentica di invocare \passthrough{\lstinline!init()!}, le funzioni
\passthrough{\lstinline!sort()!} e \passthrough{\lstinline!search()!}
lavorano su memoria sporca non inizializzata. 4. \textbf{Perdita di
Confidenzialità e Integrità}: \passthrough{\lstinline!vect!} è
accessibile in lettura/scrittura da \emph{qualsiasi} funzione del
programma. Su sistemi reali questo genera il cosiddetto
\textbf{Spaghetti Code}. - \textbf{La Regola delle Linee di Codice
(LoC)}: - Per script \textless{} 30 LoC: lo stile procedurale è rapido e
facile da mantenere. - Per progetti \textgreater{} 1000 LoC con team
concorrenti: lo stile procedurale collassa; solo l'OOP permette di
modularizzare garantendo la manutenibilità. - \textbf{La Soluzione a
Oggetti}: - Unire nello stesso modulo sia i dati
(\textbf{attributi/campi}) sia le operazioni su di essi
(\textbf{metodi}). - L'oggetto diventa l'unico responsabile della
consistenza del proprio stato.

\begin{center}\rule{0.5\linewidth}{0.5pt}\end{center}

\paragraph{C. Modellazione UML e Concetti Fondamentali di OOP (Slide
24--36)}\label{c.-modellazione-uml-e-concetti-fondamentali-di-oop-slide-2436}

\begin{itemize}
\tightlist
\item
  \textbf{UML (\emph{Unified Modeling Language})}: linguaggio grafico
  standard per specificare e documentare sistemi ad oggetti.
\item
  \textbf{Differenza Ontologica tra Classe e Oggetto}:

  \begin{itemize}
  \tightlist
  \item
    \textbf{Classe}: è il \emph{tipo} (blueprint / modello astratto).
    Definisce gli attributi e i metodi. Non occupa memoria nello heap
    per i dati dell'istanza finché non viene istanziata.
  \item
    \textbf{Oggetto}: è l'\textbf{istanza concreta} nello \textbf{Heap}.
    Possiede:

    \begin{enumerate}
    \def\labelenumi{\arabic{enumi}.}
    \tightlist
    \item
      \textbf{Stato}: i valori correnti associati ai suoi attributi.
    \item
      \textbf{Comportamento}: i metodi che può eseguire.
    \item
      \textbf{Identità}: indirizzo di memoria univoco che lo distingue
      da tutti gli altri oggetti, anche se aventi lo stesso identico
      stato!
    \end{enumerate}
  \end{itemize}
\item
  \textbf{Rappresentazione grafica UML}:

  \begin{itemize}
  \tightlist
  \item
    \textbf{Classe (Rettangolo a 3 sezioni)}:

    \begin{enumerate}
    \def\labelenumi{\arabic{enumi}.}
    \tightlist
    \item
      Nome della classe (es. \passthrough{\lstinline!Car!}).
    \item
      Attributi con tipo (es.
      \passthrough{\lstinline!manufacturer : string!},
      \passthrough{\lstinline!model : string!}).
    \item
      Metodi con segnatura (es.
      \passthrough{\lstinline!canFly() : boolean!},
      \passthrough{\lstinline!setFrameNumber(num : uint)!}).
    \end{enumerate}
  \item
    \textbf{Oggetto / Istanza (Rettangolo con nome sottolineato)}:

    \begin{itemize}
    \tightlist
    \item
      Es. \passthrough{\lstinline!Car1: Car!} con
      \passthrough{\lstinline!manufacturer = DeLorean!},
      \passthrough{\lstinline!model = DMC-12!}.
    \end{itemize}
  \end{itemize}
\item
  \textbf{Interazione a Passaggio di Messaggi (\emph{Message Passing})}:

  \begin{itemize}
  \tightlist
  \item
    Gli oggetti non leggono né modificano la memoria altrui
    direttamente: inviano richieste di servizio (\emph{messaggi}).
  \item
    L'\textbf{interfaccia} dell'oggetto è l'insieme di messaggi che esso
    accetta. Se invio un messaggio non previsto, si genera un errore di
    compilazione.
  \end{itemize}
\item
  \textbf{Visibilità}:

  \begin{itemize}
  \tightlist
  \item
    \passthrough{\lstinline!+!} (\passthrough{\lstinline!public!}):
    accessibile da qualsiasi classe.
  \item
    \passthrough{\lstinline!-!} (\passthrough{\lstinline!private!}):
    accessibile esclusivamente dall'interno della classe stessa.
  \end{itemize}
\item
  \textbf{Incapsulamento vs Information Hiding}:

  \begin{itemize}
  \tightlist
  \item
    \textbf{Incapsulamento}: impacchettare dati e codice in una singola
    unità logica (la classe).
  \item
    \textbf{Information Hiding}: nascondere i dettagli implementativi
    (rendendo i campi \passthrough{\lstinline!private!}) ed esporre solo
    un'interfaccia controllata (metodi
    \passthrough{\lstinline!public!}).
  \item
    \emph{Vantaggio chiave}: se cambiamo la struttura dati interna (es.
    da array a tabella hash), il codice esterno che usa i metodi
    pubblici non subisce alcun impatto (\emph{zero refactoring
    esterno}).
  \end{itemize}
\item
  \textbf{Associazioni UML}:

  \begin{itemize}
  \tightlist
  \item
    \textbf{Associazione standard}: linea continua con cardinalità (es.
    \passthrough{\lstinline!1!} a \passthrough{\lstinline!0..n!} tra
    \passthrough{\lstinline!Insurance!} e
    \passthrough{\lstinline!Vehicle!}).
  \item
    \textbf{Ereditarietà (Generalizzazione)}: freccia con punta
    triangolare vuota orientata verso la superclasse
    (\passthrough{\lstinline!Car!} estende
    \passthrough{\lstinline!Vehicle!}, \passthrough{\lstinline!Truck!}
    estende \passthrough{\lstinline!Vehicle!}).
  \end{itemize}
\item
  \textbf{I 4 Pilastri dell'OOP}:

  \begin{enumerate}
  \def\labelenumi{\arabic{enumi}.}
  \tightlist
  \item
    \textbf{Incapsulamento} (\emph{Encapsulation}): bundling e
    protezione dello stato interno.
  \item
    \textbf{Astrazione} (\emph{Abstraction}): separazione netta tra cosa
    un componente fa (interfaccia pubblica) e come lo fa
    (implementazione privata).
  \item
    \textbf{Ereditarietà} (\emph{Inheritance}): meccanismo di riuso e
    specializzazione gerarchica di proprietà e comportamenti.
  \item
    \textbf{Polimorfismo} (\emph{Polymorphism}): capacità di trattare
    istanze di classi diverse in modo uniforme attraverso un tipo o
    interfaccia comune.
  \end{enumerate}
\end{itemize}

\begin{center}\rule{0.5\linewidth}{0.5pt}\end{center}

\subsubsection{2. Codice: Esempio Comparativo: Dal Vettore C all'Oggetto
Java}\label{codice-esempio-comparativo-dal-vettore-c-alloggetto-java}

Per apprezzare concretamente i concetti di T02, confrontiamo il
frammento C difettoso mostrato dal docente con la sua trasposizione
nativa ad oggetti in Java.

\paragraph{Il Problema in C (Codice procedurale
insicuro)}\label{il-problema-in-c-codice-procedurale-insicuro}

\begin{lstlisting}[language=C]
// C procedurale: il dato è nudo e vulnerabile
int vect[20];
void init() { /* ... */ }

int main() {
    vect[25] = 999; // Corruzione di memoria silenziosa! Nessun errore a runtime.
    return 0;
}
\end{lstlisting}

\paragraph{La Soluzione ad Oggetti in Java (Incapsulamento, Information
Hiding,
Costruttore)}\label{la-soluzione-ad-oggetti-in-java-incapsulamento-information-hiding-costruttore}

In Java, l'oggetto controlla se stesso. Non esiste memoria non
inizializzata e l'accesso fuori limite lancia un'eccezione esplicita a
runtime:

\begin{lstlisting}[language=Java]
public class SafeVector {
    // 1. Information Hiding: il dato interno è inaccessibile all'esterno
    private final int[] elements;

    // 2. Garanzia di Inizializzazione: il costruttore viene eseguito all'atto del 'new'
    public SafeVector(int size) {
        if (size <= 0) {
            throw new IllegalArgumentException("La dimensione deve essere positiva");
        }
        this.elements = new int[size]; // Nello Heap, azzerato di default a 0
    }

    // 3. Controllo dei confini e integrità: interfaccia controllata
    public void set(int index, int value) {
        // La JVM effettua l'ArrayBoundsCheck automatico, impedendo corruzioni di memoria
        this.elements[index] = value;
    }

    public int get(int index) {
        return this.elements[index];
    }

    public int size() {
        return this.elements.length;
    }
}
\end{lstlisting}

\begin{center}\rule{0.5\linewidth}{0.5pt}\end{center}

\begin{studynote}[Punto Importante per l'Esame] \textbf{Takeaway per l'esame}: - Una classe non alloca
memoria per lo stato fino alla chiamata a \passthrough{\lstinline!new!}
(che alloca l'oggetto nello \textbf{Heap}). - L'Information Hiding si
realizza marcando i campi di istanza come
\passthrough{\lstinline!private!}. - I 4 pilastri dell'OOP
(\emph{Encapsulation}, \emph{Abstraction}, \emph{Inheritance},
\emph{Polymorphism}) sono il filo conduttore dell'intero corso.
\end{studynote}

\begin{center}\rule{0.5\linewidth}{0.5pt}\end{center}

\begin{studynote}[Nota di Studio] Il prossimo blocco è la \textbf{Lezione 3 / Slide T03:
\emph{Building, Running and Deployment in Java}}, in cui studieremo: -
Il modello di esecuzione Java: \textbf{Sorgente (.java) \(\rightarrow\)
Bytecode (.class) \(\rightarrow\) JVM / JIT}. - I comandi CLI
fondamentali: \passthrough{\lstinline!javac!},
\passthrough{\lstinline!java!}, \passthrough{\lstinline!jar!},
\passthrough{\lstinline!javadoc!}. - I concetti di \textbf{Classpath}
(\passthrough{\lstinline!-cp!}) e gestione delle dipendenze. -
\textbf{Apache Maven}: cicli di vita (\emph{compile, test, package,
install}), struttura directory standard
(\passthrough{\lstinline!src/main/java!},
\passthrough{\lstinline!src/test/java!}) e il
\passthrough{\lstinline!pom.xml!}.
\end{studynote}

\textbf{Dimmi quando sei pronto per passare a T03!}


\end{document}
